\section{Limitations and the public release}\label{sec:limits}

\subsection{What the tests do not show}
The tests measure one agent answering one question. The case the brain was built for, many agents extending one body of work over weeks, is supported by the record of Section~\ref{sec:collective} but has not been measured by a controlled comparison; the one comparison on new problems did not favour the brain, and the store held almost nothing on those problems. Section~\ref{sec:campaign} describes the comparison that will settle this.

Every test used agents of two vendors at particular model versions. Models change every few months, and a stronger model may need the store less for material it can derive and as much for material it cannot know. The measured effects are therefore statements about these agents and this store, and private question banks support controlled regression runs; they do not by themselves prevent future training-data contamination.

The no-brain arm may know published results from its training. For literature cards this makes the measured gains conservative. The store is now public, so its own cards may enter the training data of later models, and a later no-brain arm may then know the agents' own results as well; a test of a later model against this store will understate the brain's effect for a reason that has nothing to do with the brain. Fresh evaluation banks should be drawn from unseen material and checked against the released files and the model's documented training boundary where available.

\subsection{What the store does not guarantee}
Most cards are \textsf{unverified} or \textsf{provisional}. The store records what was claimed, by what route and with what support; it does not certify that the claims are true. The curated secondary cards were written by agents from textbooks and surveys, and a card can misstate a hypothesis of its source; the fidelity of statement to source has been reviewed for many cards but not all. A reader who uses a card in a proof should read its source.

The store is unevenly distributed. It is dense in probability theory, Boolean functions and the theory of deep learning, where the campaigns were, and thin elsewhere; $2{,}878$ cards have no link and can be found only by their words.

The retrieval is a word search with a meaning channel. It does well when the question can be worded the way a card is, and poorly when it cannot: the hook was silent on all $60$ reworded questions of the store test. A whole-card embedding that does better on reworded questions has been built and not yet switched on.

\subsection{Cost}
A search prints about $11{,}000$ characters, and an agent that searches many times pays for every printout again at each later turn of its session. The measured savings of Section~\ref{sec:tests} come from walking instead of searching again; an agent that ignores the advice of Section~\ref{sec:usage} can spend more with the brain than without it. On the repository test the brain arm read $14\%$ more tokens per question for a gain that was not significant.

\subsection{Open access and the released snapshot}\label{sec:release}
The release is available to everyone at \url{https://github.com/yspennstate/math-brain-public}. The licence decision is settled. The original code, documentation and project-owned data are under MIT, which permits use, copying, modification and redistribution, including commercial use, with its copyright and licence notice retained. This paper is additionally made available under Creative Commons Attribution 4.0 International (CC BY 4.0), permitting sharing and adaptation with attribution. No payment, registration, institutional affiliation or individual permission is required to read or download the release.

Third-party rights are preserved. The $1{,}766$ Wikipedia-derived cards retain CC BY-SA 4.0 and the $378$ ProofWiki-derived cards retain CC BY-SA 3.0, with their source attribution and ShareAlike requirements. These are open licences that permit commercial reuse; MIT does not replace them. A citation to a book or paper, or the absence of a \texttt{source\_id}, does not establish ownership or permission to relicense its expression. The project licences cover the author's original contributions, not rights belonging to other authors. \texttt{RIGHTS.md} states this boundary.

The checked public snapshot is the export at commit \href{https://github.com/yspennstate/math-brain-public/commit/393821a830f0979d27c38430130900ac57fb1701}{\texttt{393821a830f0}}, read on 26 October 2026. Its manifest records $13{,}344$ cards, $19{,}398$ links, $231$ retractions and $1{,}052$ bundle manifests, in $56{,}733$ rows across sixteen tables. All sixteen file checksums and row counts matched, and a clean import produced no foreign-key violations and no card or bundle-manifest digest failures. This supersedes the obsolete 12--13 October export description. The dated 22 October census elsewhere in the paper remains a historical measurement, not a count of this later snapshot. The release has its own Git history; it is not a publication of the private working repositories' histories.

The export's portable digests verify its released contents. They do not certify source fidelity, copyright clearance, mathematical correctness, or completion of a new-revision repair in the private store. Historical campaign manifests may hash artifacts before sanitization. The retained papers concern mathematics, mathematical physics and the mathematical-research system; this description is not a certification of every file copied into their supporting directories.

Public availability and completion of every audit are different facts. The release revision verifies the exported counts and digests and corrects the licensing statement. It does not assert that every book-derived passage has been compared with its source, that every private evaluation bank has been checked against this enlarged snapshot, or that all historical files have passed a single comprehensive privacy review. Long third-party quotations require a source-specific rights review and, where necessary, shortening or removal. Items answerable from the public release must be excluded from fresh held-out claims; repetitions on those items can be reported only as regression tests. These limits remain explicit rather than being converted into unsupported completion claims.
